\documentclass[12pt,a4paper,twoside]{book}
\usepackage[left=1.00in, right=1.00in, top=1.00in, bottom=1.00in]{geometry}
\usepackage[spanish,mexico]{babel}
\usepackage{graphicx}
\graphicspath{{img/}}  % Incluir la ruta de la carpeta de imagenes
\usepackage{graphics} %Incluir imagenes pdf
\usepackage{xcolor}
\definecolor{coolblack}{rgb}{0.0, 0.18, 0.39}
%%%%%%%%%%%%%%%%%%%%%%%%%    Bibliography   %%%%%%%%%%%%%%%%%%%%%%%%%%%%%%
\usepackage{csquotes}
\DeclareUnicodeCharacter{0301}{\'{e}} 
\usepackage[style=iso-authoryear]{biblatex}
\addbibresource{lib/bibliography.bib}
%%%%%%%%%%%%%%%%%%%%%%%%%%%%%%%%%%%%%%%%%%%%%%%%%%%%%%%%%%%%%%%%%%%%%%%%%%%
\usepackage{hyperref} % Carga hyperref primero
\hypersetup{
    colorlinks=true,
    citecolor=coolblack, %color citas
    linkcolor=black, %Cambiar al color de preferencia
    filecolor=magenta,      
    urlcolor=cyan,
    pdftitle={Overleaf Example},
    % pdfpagemode=FullScreen,
}
\usepackage[
type={CC},
modifier={by-nc-nd},
version={4.0},
]{doclicense}
\usepackage{url} 
\urlstyle{same}
\usepackage{hyperxmp} % Luego carga hyperxmp
%%%%%%%%%%%%%%%%%%%%%%%%%%%%%%%%%%%%%%%%%%%%%%%%%%%%%%%%%%%%%%%%%%%%%%%%%%%
%%%%%%%%%%%%%%%%%%%%%%%%%% Change font %%%%%%%%%%%%%%%%%%%%%%%%%%%%%%%%%%%%
% Configuración adicional
%\renewcommand{\familydefault}{\sfdefault} % Cambiar la fuente por defecto a sans-serif
\usepackage[T1]{fontenc}
\usepackage{setspace}
%\singlespacing  % Espaciado simple
\onehalfspacing % Espaciado 1.5
% \doublespacing

\usepackage{newtxtext,newtxmath} % Paquete para usar Times New Roman
%\usepackage{cmbright} % Opcional: Computer Modern Bright
%\usepackage{lmodern} % Opcional: Latin Modern
%\usepackage{times}    % Opcional: Times New Roman
% \usepackage{tgtermes}
%\usepackage{helvet}   % Opcional: Arial
%%%%%%%%%%%%%%%%%%%%%%%%% Chapter Configuration %%%%%%%%%%%%%%%%%%%%%%%%%%%
% Paquete para configurar títulos
\usepackage{titlesec}

% Configuración de numeración
\renewcommand{\thechapter}{\Roman{chapter}}
\renewcommand{\thesection}{\arabic{chapter}.\arabic{section}}
\renewcommand{\thesubsection}{\thesection.\arabic{subsection}}
\renewcommand{\thesubsubsection}{\thesubsection.\arabic{subsubsection}}
\renewcommand{\theparagraph}{\thesubsubsection.\arabic{paragraph}}
\renewcommand{\thesubparagraph}{\theparagraph.\arabic{subparagraph}}

% Configuración de tamaños de títulos
\titleformat{\chapter}[display]
{\normalfont\bfseries\centering}{\MakeUppercase{\chaptertitlename}\ \thechapter}{0pt}{\normalsize}
\titleformat{\section}
{\normalfont\bfseries}{\thesection}{1em}{\normalsize}
\titleformat{\subsection}
{\normalfont\bfseries}{\thesubsection}{1em}{\normalsize}
\titleformat{\subsubsection}
{\normalfont\bfseries}{\thesubsubsection}{1em}{\normalsize}
\titleformat{\paragraph}
{\normalfont\bfseries}{\theparagraph}{1em}{\normalsize}
\titleformat{\subparagraph}
{\normalfont\bfseries}{\thesubparagraph}{1em}{\normalsize}

% Espaciado de los títulos
\titlespacing*{\chapter}{0pt}{0pt}{20pt}
\titlespacing*{\section}{0pt}{12pt}{6pt}
\titlespacing*{\subsection}{0pt}{12pt}{6pt}
\titlespacing*{\subsubsection}{0pt}{12pt}{6pt}
\titlespacing*{\paragraph}{0pt}{12pt}{6pt}
\titlespacing*{\subparagraph}{0pt}{12pt}{6pt}

% Configuración de numeración hasta el cuarto nivel
\setcounter{secnumdepth}{3}
\setcounter{tocdepth}{3}

%%%%%%%%%%%%%%%%%%%%%%%%% Configuration Table of contents %%%%%%%%%%%%%%%%%
\usepackage{tocloft}
\renewcommand{\cfttoctitlefont}{\hfill\normalsize\bfseries\MakeUppercase} % Cambia el tamaño y estilo de "Índice General"
\renewcommand{\cftaftertoctitle}{\hfill} % Centra el título "Índice General"
\renewcommand{\cftlottitlefont}{\hfill\normalsize\bfseries\MakeUppercase} % Cambia el tamaño y estilo de "Índice de Tablas"
\renewcommand{\cftafterlottitle}{\hfill} % Centra el título "Índice de Tablas"
\renewcommand{\cftloftitlefont}{\hfill\normalsize\bfseries\MakeUppercase} % Cambia el tamaño y estilo de "Índice de Figuras"
\renewcommand{\cftafterloftitle}{\hfill} % Centra el título "Índice de Figuras"

\usepackage{titletoc}
\titlecontents{chapter}[1.5em]
{\vspace{10pt}\bfseries} % Formato del título
{\MakeUppercase{CAPÍTULO} \thecontentslabel {\titlerule*[0.7pc]{.} \contentspage} \\ } % Etiqueta (número de capítulo) Aquí se configura \\
{} % Espacio entre la etiqueta y el título
{\titlerule*[0.7pc]{.}\contentspage} % Rellenado y número de página

\titlecontents{section}[3.8em]
{} % Sin formato del título
{\thecontentslabel\hspace{1em}} % Etiqueta (número de sección)
{} % Espacio entre la etiqueta y el título
{\titlerule*[0.7pc]{.}\contentspage} % Rellenado y número de página

\renewcommand{\cftchapfont}{\normalfont\bfseries}
\renewcommand{\cftsecnumwidth}{3em} % Ajusta el ancho del número de sección
\renewcommand{\cftsubsecnumwidth}{3em} % Ajusta el ancho del número de subsección
\setlength{\cftchapnumwidth}{4em} % Ajusta el espacio para el número de capítulo
%%%%%%%%%%%%%%%%%%%%%%%%%%%%%%%%%%%%%%%%%%%%%%%%%%%%%%%%%%%%%%%%%%%%%%%%%%%
\usepackage{lastpage}
\usepackage{fancyhdr}
\pagestyle{fancy}
\fancyhf{} % Limpia todos los encabezados y pies de página
\fancyhead[R]{- \thepage\ de \pageref{LastPage} -} % Configura el encabezado derecho

% Asegúrate de que el estilo de página 'plain' también use el mismo encabezado
\fancypagestyle{plain}{
    \fancyhf{}
    \fancyhead[R]{- \thepage\ de \pageref{LastPage} -}
}
\renewcommand{\headrulewidth}{0pt}
\usepackage{appendix}
\usepackage{acronym}
\usepackage{mdframed}
\usepackage{lipsum}
% \usepackage{showframe}
\author{Juan Benito Quintana Arone}
\title{"Tesis 2024"}
%%%%%%%%%%%%%%%%%%%%%%%%%%%%%%%%%%%%%%%%%%%%%%%%%%%%%%%%%%%%%%%%%%%%%%%%%%%
\begin{document}
\frontmatter
    
\begin{titlepage}

	\begin{center}
		\textbf{\large{UNIVERSIDAD NACIONAL MICAELA BASTIDAS DE APURÍMAC}}\\
		\vspace{2mm}
		\large{\textbf{Faculdad de Ingeniería}}\\
		\vspace{2mm}
		\large{\textbf{}{Escuela académico profesional de Ingeniería Civil}}\\
		%\vspace{2mm}
		%\normalsize{\textbf{SUB SEDE ABANCAY}}
		\vspace{5mm}
		\begin{figure}[h!]
			\centering
			\includegraphics[width=0.3\linewidth]{0.2.pdf}
		\end{figure}
		
		\vspace{5mm}
		\large{Proyecto de Tesis:}\\
		%\vspace{5mm}
		Vulnerabilidad y peligro sísmico en viviendas construidas con adobe en la Comunidad Campesina de Paccaypata\\
		\vspace{10mm}
		Presentado por:\\
		\vspace{10mm}
		Juan Benito Quintana Arone\\
		\vspace{5mm}
		Para optar el título profesional de Ingeniero Civil 
		
		\vspace{10mm}
		
		\vfill
		Abancay, Perú\\
		\vspace{5mm}
		2024\\
	\end{center}
\end{titlepage}


    \begin{titlepage}
	\begin{center}
		\textbf{\large{UNIVERSIDAD NACIONAL MICAELA BASTIDAS DE APURÍMAC}}\\
		\textbf{\large{FACULTAD DE INGENIERÍA}\\}
		\large{ESCUELA ACADÉMICO PROFESIONAL DE INGENIERÍA CIVIL}\\
		\vspace{2mm}
		\begin{figure}[h!]
			\centering
			\includegraphics[width=0.2\linewidth]{0.1.pdf}
		\end{figure}
		%\vspace{2mm}	
		\textbf{TESIS}\\
		\vspace{2mm} 
		\uppercase{\textbf{Impacto de los cambios en los costos de proyectos de inversión pública del gobierno regional de apurímac, periodo 2018 - 2022}}\\
		\vspace{3mm}
		Presentado por \textbf{Juan Benito Quintana Arone}, para optar el titulo de: Ingeniero Civil \\
		\vspace{2mm}

		\begin{center}
			\textbf{Presidente:} \hfill \rule{9cm}{0.4pt}\\
			\hspace{7cm} \textnormal{\sffamily Ph.D: Lucy M. Guanuchi Orellana} \\
			\vspace{5mm}
			\textbf{Primer miembro	:}\hfill \rule{9cm}{0.4pt}\\
			\hspace{7cm} \textnormal{\sffamily Dr: Nelson Palemon Meza Peña} \\		
			\vspace{5mm}
			\textbf{Segundo miembro	:}\hfill \rule{9cm}{0.4pt}\\
			\hspace{7cm} \textnormal{\sffamily Mstro: Feliciano Escobedo Silva} \\		
			\vspace{5mm}
			\textbf{ \large ASESOR	:} \hfill \rule{9cm}{0.4pt}\\	
			\hspace{7cm} {\normalsize Ing. Mgt. Diómdes N. Ferrel Sarmiento} \\	
		\end{center}
		
		\vfill
		
	\end{center}
\end{titlepage}
    % \newpage
\pagestyle{plain}
\vspace*{\fill} %Centrar verticalmente en la pagina
\begin{flushright}%Alinea a la izquierda
	\begin{minipage}[t]{0.8\textwidth}
		{\large\textbf{Agradecimientos}\\}
		En primer lugar, quiero agradecer a Dios por su amor y guía constante en mi vida. Ha sido mi fuerza y mi refugio en los momentos de desafío. Su bendición ha sido fundamental en cada paso del camino.
		
		A mis padres Teresa y Benito, quiero expresarles mi profundo agradecimiento por su amor inquebrantable y su apoyo incondicional. Han sido mis pilares, siempre presentes para brindarme palabras de aliento, ánimo y sabiduría en cada etapa de mi vida universitaria. 
		
		A mis hermanos, gracias por ser mis compañeros de vida y mis mayores motivadores. Su apoyo incondicional y su fe en mí me han impulsado a seguir adelante incluso en los momentos más desafiantes. 
		
		\vspace{5mm}
		
		Con profundo amor y gratitud,\\
		\vspace{5mm}
		Juan Benito Quintana Arone
		
		
		
	\end{minipage}
\end{flushright}

\vspace*{\fill}
    % \newpage
\pagestyle{plain}
\vspace*{\fill} %Centrar verticalmente en la pagina
\begin{flushright}%Alinea a la izquierda
	\begin{minipage}[t]{0.8\textwidth}
		{\Large\textbf{Dedicatoria}\\}
		
	En este día tan especial, quiero dedicarles mi tesis a mis padres Teresa, Benito y mis hermanos como una muestra de gratitud y reconocimiento a cada uno de ustedes. Su presencia en mi vida ha sido un regalo inmenso y estoy infinitamente agradecido por tenerlos a mi lado.
		
	\end{minipage}
\end{flushright}

\vspace*{\fill}
    \newpage
\pagestyle{plain}
%\vspace*{20cm}
%\vspace*{\fill}
\noindent Impacto de los cambios en los costos de proyectos de inversión pública del gobierno regional de apurímac, periodo 2018 - 2022\

\noindent Linea de investigación denominado ''Ingeniería de la conctrucción''
%El autor autoriza a la Universidad Nacional Micaela Bastidas de Apurímac a reproducir la tesis en su totalidad o en parte, con fines estrictamente académicos.
%En particular esta licencia no permite venta ni modificaciones de este material.\\
\doclicenseThis%

	




    \tableofcontents
    \newpage
	\listoftables
	\newpage
	\listoffigures
	\newpage
    \chapter*{SIGLAS Y ACRÓNIMOS}

\begin{acronym}
	% \acro{pdf}{Portable Document Format}
	% \acro{html}{Hypertext Markup Language}
	\acro{IGP}{Instituto geofísico del Perú}
	\acro{INDECI}{Instituto Nacional de Defensa Civil}
	\acro{Ingemment}{ Instituto Geológico Minero y Metalúrgico}
	\acro{INEI}{Instituto Nacional de Estadística e Informática}

\end{acronym}









\mainmatter     
    \pagestyle{headings}
\clearpage
\chapter*{INTRODUCCIÓN}
La medición del Índice de Rugosidad Internacional (IRI) es un método utilizado para evaluar la calidad de las carreteras y la comodidad de los usuarios al transitar por ellas. El IRI se define como la medida del movimiento vertical de un vehículo que se produce debido a las irregularidades de la superficie de la carretera. Esta medida se expresa en unidades de longitud, generalmente en metros por kilómetro (m/km)o también (mm/m).

La medición del IRI se realiza mediante el uso de equipos especiales que registran el movimiento vertical del vehículo a lo largo de la carretera. Estos equipos están equipados con acelerómetros y otros sensores que miden la aceleración y la velocidad del vehículo en tiempo real. La información obtenida se procesa posteriormente mediante un software especializado que calcula el valor del IRI.

El IRI se utiliza como un indicador de la calidad de la superficie de la carretera y su influencia en el confort y la seguridad de los usuarios. Un IRI alto indica una carretera con superficie irregular y puede provocar fatiga del conductor, aumentar el consumo de combustible, generar daños en los vehículos y disminuir la vida útil de la carretera. Por lo tanto, la medición del IRI es esencial para la planificación y mantenimiento adecuado de las carreteras.

%La \acf{AI} es una rama de la ciencia de la computación que busca desarrollar sistemas que imiten la inteligencia humana. El \ac{ML} y el \acl{DL} son subcampos importantes dentro de la \ac{AI}. Además, el \acs{NLP} es otra área de investigación relevante que se enfoca en la interacción entre humanos y máquinas a través del lenguaje.


%\Gls{latex} es muy popular en el ámbito académico.\acrfull{pdf}





    \newpage
\chapter*{Datos generales}

\section*{Título del proyecto}
Impacto de los cambios en los costos de proyectos de inversión pública del Gobierno Regional de Apurímac, periodo 2018 - 2022
\section*{Ejecutor} 
\noindent Juan Benito Quintana Arone

\section*{Asesor y Coasesor}
\noindent Mgt.Diómedes Napoleón Ferrel Sarmiento , Asesor\\
Nombres y Apellidos (en minúsculas), coasesor

\section*{Línea de investigación}
\noindent Ingeniería de la construcción

\newpage


    \chapter*{RESUMEN}

En resumen 


\lipsum[5] 
    
\chapter*{ABSTRACT}
	
	\lipsum[2-4]


    \pagestyle{fancy}
\chapter{PLANTEAMIENTO DEL PROBLEMA}
\section{Descripción del problema}
\lipsum[4-5]
\section{Enunciado del problema}
\subsection{Problema general}
\subsection{Problemas específicos}
\section{Justificación de la investigación}

    \chapter{OBJETIVOS E HIPÓTESIS}
\section{Objetivos de la investigación}
\subsection{Objetivo general}
\subsection{Objetivos específicos}
\section{Hipótesis de la investigación}
\subsection{Hipótesis general}
\subsection{Hipótesis específicas}
\section{Operacionalización de variables}

    \chapter{MARCO TEÓRICO REFERENCIAL}
\section{Antecedentes}
% Necesitas saber cuánto se ha estudiado las variables
% Antecedentes mundiales
\cite{JimenezRios2023} realizó un estudio titulado: Learning from the Past: Parametric Analysis of Cob Walls en la ciudad Oslo del país Irlanda.El objetivo general del estudio fue explorar las ventajas y desventajas de cada uno de ellos y proporcionar alternativas a los profesionales: análisis FEM elástico lineal y análisis límite cinemático con macro-elementos. La muestra estuvo constituida por muros de adobe o "cob" como ellos de denominan en Irlanda. El diseño que se utilizó fue experimental. Los instrumentos que se usaron fueron el modelado FEM,  siguiendo dos enfoques computacionales, a saber, el método de elementos finitos y el análisis límite cinemático.Y los resultados obtenidos demostraron que los muros tradicionales de barro en Irlanda son muy robustos. Se requerirían valores de aceleración relativamente altos, poco probables en una región con bajo riesgo sísmico como Irlanda.

\cite{Wang2023} realizó un estudio titulado: Adobe-Brick-Masonry Composite Wall with a Wooden-Construction Center Column en la ciudad Shihezi del país China.El objetivo general del estudio fue proponer una tecnología para mejorar el desempeño sísmico de un muro compuesto de albañilería de adobe modificado con una columna central de construcción en madera La muestra estuvo constituida por muro de adobe modificado.El diseño que se utilizó fue experimental.Los instrumentos que se usaron fueron  pruebas de carga cíclica quasiestáticas en plano para estudiar sus modos de fallo e indicadores de desempeño sísmico en un software y los resultados obtenidos han sido mejoraron su capacidad de carga sísmica gracias al barro modificado por otra parte la columna central y los materiales de amarre proporcionaron una segunda línea de defensa que aumentó la ductilidad de los muros y la zona residual de colapso.

\cite{Romanazzi2023} realizó un estudio titulado:  Effectiveness of a TRM solution for rammed earth under in-plane cyclic loads en la ciudad Braga y Guimarães del país Portugal.El objetivo general del estudio fue evaluar la efectividad de una solución de refuerzo TRM para paredes de tierra apisonada sometidas a cargas cíclicas en el plano.una pared de tierra apisonada de un edificio tradicional de un solo piso con techo de madera y geometría en forma de I en planta. El diseño que se utilizó fue experimental.Los instrumentos que se usó fué protocolo de prueba del modelo reforzado GeoRE-IP Y los resultados obtenidos destacaron la efectividad de la solución TRM en la mejora de la capacidad de corte en el plano, la ductilidad y la energía disipada del modelo.

% Antecedentes latinoamericanos
\cite{LopezP.2020} realizó un estudio titulado: Comportamiento sísmico de edificaciones de tapia pisada reforzadas con marcos de madera y viga de coronación en concreto, en la ciudad Bogotá del país Colombia.El objetivo general del estudio fue reducir esta vulnerabilidad.La muestra estuvo constituida por  ensayos en mesa vibratoria en un modelo a escala 1:6 y su contraparte no reforzada .El diseño que se utilizó fue experimental.Los instrumentos que se usaron fueron Implementación de un diafragma rígido y los resultados obtenidos han sido  muestran una mejora significativa en el comportamiento global de la construcción (reducción en los desplazamientos hasta del 80\% en los modelos reforzados), logrando el objetivo de prevenir el colapso, manteniendo la integridad de la estructura.

\cite{TorresMoreno2022} realizó un estudio titulado: Influencia de la resistencia a la compresión y módulo de elasticidad del adobe en el comportamiento estructural de viviendas patrimoniales en la ciudad Riobamba del país Ecuador.El objetivo general del estudio fue obtener una base de datos de la resistencia a la compresión y módulo de elasticidad del adobe, y determinar su influencia en el comportamiento estructural de tres viviendas patrimoniales en el cantón Guamote. La muestra estuvo constituida por 12 muestras de adobe de 3 viviendas diferentes El diseño que se utilizó fue experimental-correlacional. Los instrumentos que se usaron fueron ensayos en el laboratorio  los resultados obtenidos muestran que las resistencias a la compresión que van de 0.062 a 0.1852 MPa y módulos de elasticidad de 0.9012 a 2.1195MPa. 

\cite{AmaguayBermeo2022} realizó un estudio titulado: Evaluación y reforzamiento estructural, incorporando mampostería enchapada y alternativa para el mejoramiento de suelo o refuerzo de cimentación de una edificación que presenta asentamientos diferenciales en la ciudad de Quito del país Ecuador.El objetivo general del estudio fue evaluación estructural a una edificación de departamentos ubicada en la parroquia de la Magdalena.La muestra estuvo constituida por una sola vivienda de mampostería.El diseño que se utilizó fue no experimental Los instrumentos que se usaron fueron ficha de recolección de datos, mediante una inspección visual y los resultados obtenidos han sido que la vivienda en estudio presenta una alta vulnerabilidad e implementó un reforzamiento enchapado.

% Antecedentes nacionales y regionales
\cite{OteroMonteza2022} realizó un estudio titulado: Influence of the aspect ratio on seismic performance of adobe buildings en la ciudad de Lima del país Perú. El objetivo general del estudio fue el comportamiento sísmico de construcciones de adobe con muros más grandes en un eje que en el otro eje ortogonal. La muestra estuvo constituida por cuatro modelos de edificios con diferentes relaciones de aspecto, que varían en longitud desde los ocho metros hasta los cincuenta y dos metros.  El diseño que se utilizó fue experimental Los instrumentos que se usaron fueron análisis de historia del tiempo no lineal con tres registros sísmicos peruanos Y los resultados obtenidos han sido Muestra los modos fundamentales de vibración en la dirección del eje y el primer modo de vibración involucra el 48.48\% de la masa total.

\cite{QuirozHuaraya2021} realizó un estudio titulado: Propuesta de reforzamiento estructural de viviendas de adobe utilizando mallas de acero electrosoldadas en Huarangal, Carabayllo en la ciudad de Lima del país Perú.El objetivo general del estudio fue proponer un tipo de reforzamiento para muros de adobe con el uso de mallas electrosoldadas.La muestra estuvo constituida En la primera fase se realizaron ensayos de compresión axial en tres pilas de adobe y compresión diagonal en tres muretes de adobe sin reforzamiento. En la segunda fase se ensayaron 12 muretes de adobe reforzados con dos tipos de mallas electrosoldadas y en diferentes lados de los especímenes. El diseño que se utilizó fue experimental.Los instrumentos que se usaron fueron ensayo de mueretes con mallas electrosoldadas y sin mallas y los resultados obtenidos han sido los ensayos demostraron un incremento en la resistencia al corte y una mejora en la ductilidad de los muretes de adobe reforzados debido al tipo de malla y al número de lados reforzados en cada murete.

\cite{CerveraTimana2023} realizó un estudio titulado: Evaluación de la vulnerabilidad sísmica en las edificaciones de la zona sur-este del distrito de Lambayeque en la ciudad de Chiclayo del país Perú.El objetivo general del estudio fue determinar los niveles de vulnerabilidad sísmica en las edificaciones existentes de la zona sur-este del distrito de Lambayeque. La muestra estuvo constituida por se evaluaron 3,054 edificaciones.El diseño que se utilizó fue no experimental Los instrumentos que se usaron fueron Método del Índice de Vulnerabilidad de Benedetti - Petrini y los resultados obtenidos muestran que 477 edificaciones resultaron con baja vulnerabilidad representando un 15.62\%; 1,901 edificaciones resultaron con vulnerabilidad media representando un 62.25\%; y 676 edificaciones resultaron con alta vulnerabilidad representando un 22.13\%.

\section{Marco teórico}

\cite{CENEPRED2014}

\section{Marco conceptual}
    \chapter{METODOLOGÍA}
\section{Tipo y nivel de investigación}
\subsection{Tipo de la investigación}

Se presentan diversas perspectivas con respecto a la clasificación o división de la investigación. No obstante, independientemente del criterio utilizado, según \cite[]{Zacarias2020} se espera que cumpla con el principio de parsimonia, es decir, que sea exhaustivo y excluyente.

\cite[47]{SanchezCarlessi2015} destaca la importancia de estar conscientes sobre la naturaleza y los propósitos de la investigación, clasificándolos en tres categorías. En primer lugar, se encuentra la investigación básica, también conocida como pura o fundamental, según lo menciona el autor. En este tipo de investigación, el investigador sostiene que \emph{"[...] Mantiene como propósito recoger información de la realidad para enriquecer el conocimiento científico, está orientado al descubrimiento de principios y leyes”}. En segunda instancia, nos encontramos con la investigación aplicada, también denominada constructiva o utilitaria. En este caso, la investigación aplica los conocimientos teóricos a una situación específica para obtener consecuencias prácticas derivadas de la deducción. Por último, se aborda la investigación sustantiva, donde el enfoque está dirigido a describir, explicar y prever la realidad.

Por otra parte, \cite[36]{Tamayo2004} clasifica la investigación según su enfoque. Por un lado, se encuentra la investigación pura, también denominada básica o fundamental. Esta modalidad se caracteriza por buscar confrontar la teoría con la realidad, con la finalidad de desarrollar una teoría basada en leyes y principios. Por otro lado, está la investigación aplicada, también conocida como activa o dinámica. Su particularidad radica en que se centra en la aplicación inmediata de conocimientos, sin enfocarse en el desarrollo de teorías.

Ahora bien, al considerar la revisión bibliográfica de los dos autores, se determina que la presente investigación adopta un enfoque básico, puro o fundamental. En este contexto, se llevará a cabo la recopilación de datos existentes en el Gobierno Regional de Apurímac.

\subsection{Nivel de la investigación} \label{nivel_invest}

\cite[88]{RiosRamirez2017} clasifica el nivel de investigación según los propósitos o finalidades, la naturaleza de los datos o enfoque, el nivel de conocimiento, las fuentes de información, las condiciones de información, el tiempo y el diseño o control de las variables. En este caso, es necesario profundizar en el criterio del nivel de conocimiento, según el autor, quien afirma que "[...] involucra el grado de conocimiento sobre el objeto de estudio". Al mismo identifica cuatro tipos de investigación según esta clasificación. El primero es el exploratorio, un estudio en el cual no se puede distinguir claramente entre las variables independientes o dependientes, de ahi que también se le conoce como univariado. En este tipo de investigación, se examinan temas poco abordados. El descriptivo, por su parte, "[...] busca encontrar las características, comportamiento y propiedades del objeto de estudio" en el tiempo presente o futuro. La investigación relacional hace mérito a su nombre, mide la relación entre dos o más variables dadas, pero es importante señalar que no determina la causa-efecto. Por último, la investigación explicativa sí determina la relación causa-efecto.

Por otra parte, si bien es cierto que \cite[93]{HernandezSampieri2014} coincide con \citeauthor{RiosRamirez2017} en la mayoría de las denominaciones de su clasificación, el autor utiliza la denominación "correlacional" en lugar de "relacional". Muchos autores como \cite[]{Zacarias2020} opinan que debería llamarse "relacional", ya que la relación implica una conexión entre variables, mientras que la correlación se refiere a la relación entre unidades.

\section{Diseño de la investigación}
\cite[128]{HernandezSampieri2014} explica el diseño de la investigación como el "plan o estrategia que se desarrolla para obtener la información requerida en una investigación". El autor distingue dos diseños, ambos con una importancia destacada. 

En primer lugar, está el diseño experimental, donde las variables independientes se manipulan deliberadamente para observar sus efectos sobre otras variables dependientes en una situación controlada. A su vez, este diseño se subdivide en tres clases. La característica distintiva de todos ellos es que deben tener validez interna y externa, es decir, que los resultados deben ser aplicables primero a la población en estudio y luego al entorno externo. Las tres subclases son: preexperimentos, experimentos puros y cuasiexperimentos. No profundizaremos en estos enfoques ya que no es el objetivo de este estudio.

En contraste, los diseños no experimentales son aquellos estudios en los cuales solo se observan los fenómenos, y no se realiza una manipulación deliberada de variables. Este tipo de diseño, según su dimensión temporal o puntos en los cuales se recolectan datos, se puede clasificar en investigación transeccional, también conocida como transversal, la cual se caracteriza por recopilar datos una sola vez y en un solo momento. A su vez, se clasifica en tres subcategorías: exploratorios, descriptivos y correlacionales-causales, basándose en lo mencionado en la sección \ref{nivel_invest}. Por otro lado, la investigación longitudinal o evolutiva obtiene datos en diferentes puntos en el tiempo, ya que compara los resultados a través del cambio.\cite[341]{NaupasPaitan2014} clasifica este diseño en descriptiva simple, descriptiva comparativa, causal-comparativa, correlacional, longitudinal y transversal.

En cuanto a los criterios de planificación para la toma de datos en el tiempo, no se encontraron autores que mencionen explícitamente esta clasificación. Sin embargo, se sobreentiende que si los datos fueron planificados previamente, el estudio será prospectivo. En cambio, si los datos ya estaban registrados, el estudio será retrospectivo.

\section{Descripción ética de la investigación}

Según \cite[]{AbreuSuarez2017}, la ética se presenta como un campo que abarca los valores esenciales del ser humano, como la honestidad, la solidaridad, el respeto y la tolerancia, entre otros aspectos. Estos elementos deben ser considerados por los investigadores.

El investigador aseguró obtener el consentimiento informado de todas las partes involucradas en la investigación, abarcando tanto a los funcionarios del Gobierno Regional de Apurímac como a cualquier otra persona que pudiera ser afectada directa o indirectamente por los resultados del estudio.

Se garantizó la confidencialidad de los datos y la privacidad de los participantes en el estudio. La información recopilada fue manejada de manera segura y exclusivamente utilizada con fines investigativos específicos, preservando la identidad y la información personal de los involucrados.

Manteniendo la imparcialidad y objetividad en todas las etapas del estudio, el investigador basó los resultados en el análisis de datos precisos y fiables, evitando sesgos y asegurando una interpretación justa y equitativa de los hallazgos.

La investigación se llevó a cabo con honestidad y rigor académico, evitando la fabricación, falsificación o manipulación de datos, así como cualquier forma de plagio de trabajos previos.

Considerando el equilibrio entre beneficios potenciales y riesgos asociados con el estudio, se garantizó que los posibles beneficios superaran cualquier daño potencial a los participantes o a la comunidad en general.

Dada la naturaleza que involucra seres humanos o datos sensibles, el investigador obtuvo la aprobación de la Subgerencia de Obras del Gobierno Regional de Apurímac, entidad encargada de supervisar la integridad y el bienestar de los participantes.

El investigador practicó la transparencia en la divulgación de métodos, resultados y conclusiones, incluyendo la revelación de cualquier conflicto de intereses y proporcionando los datos subyacentes para permitir la validación y replicación del estudio por parte de otros.

Los resultados de la investigación serán utilizados de manera responsable y ética, considerando sus implicaciones y su posible impacto en la toma de decisiones y políticas públicas.

Se mostró respeto y valoración hacia la cultura, tradiciones y conocimientos locales, evitando cualquier forma de apropiación indebida de la información recopilada.

En la medida de lo posible, la investigación buscó contribuir al bienestar social y al desarrollo de la comunidad, proporcionando información relevante y útil para mejorar la gestión de proyectos de inversión pública en el Gobierno Regional de Apurímac.

El presente trabajo de investigación en cuanto a las citaciones de referencias bibliográficas sigue las directrices de las normas establecidas por \cite[]{IOS2021}, mas conocidos como ISO-690 , en cuanto a la redacción siguen las pautas proporcionadas por el formato del Vicerrectorado de investigación de la Universidad Nacional Micaela Bastidas de Apurímac.

Por otra parte se tomó en cuenta con lo dispuesto por el \cite[]{VRI2018}, los titulos que corresponden a las normas de comportamiento de quienes investiguen en la página 2, buenas prácticas de los investigadores y la investigación con personas.

    % \input{lib/5.0.chapter/admin_project.tex}
    \input{lib/6.0.chapter/result_disc.tex}
    \input{lib/7.0.chapter/conclu_recomen.tex}
\backmatter

\lipsum[1-8]
\end{document}